% =================================================================
% Part XVI: NRMO v7.3 — Production Contract Integration
%   Reconciles the v7.3 implementation-target requirements (as
%   evidenced by DecisionCompass secondary sources, 2026-09) against
%   the v7.2 / v7.2.1 / Normative Canon baseline. Does not amend the
%   Normative Canon. See NRMOIntegrated/docs/V7_3_DIFFERENTIAL_TABLE.md
%   for the full row-by-row reconciliation this Part summarises.
% =================================================================

\part{NRMO v7.3 --- Production Contract Integration}
\label{part:v73_production_contract}

\chapter{Provenance and Status of This Part}
\label{ch:v73_provenance}

\nrmobox{Provenance and authority status --- read before use}{
\textbf{This Part is not a transcription of an original ``NRMO SYSTEM
v7.3'' specification document.} DecisionCompass Issue \#120 records
that such a document (\texttt{NRMO SYSTEM\_7.3.md} and
\texttt{NRMO\_SYSTEM\_v7\_3\_PATCH.md}) was supplied 2026-09-15, but as
of this build that original text could not be located in
\texttt{zarame96/NRMO}, \texttt{zarame96/DecisionCompass} (any branch),
or local session storage. \\[0.5em]
What this Part \emph{is}: a reconciliation of (a) the NRMO v7.2
Normative Canon, (b) the v7.2.1 Passive Ruin revision, and (c) the
secondary implementation-target documents
\texttt{docs/NRMO\_V7\_3\_IMPLEMENTATION\_SCOPE.md} and
\texttt{docs/NRMO\_V7\_3\_GAP\_MATRIX.md} from the DecisionCompass
branch \texttt{agent/nrmo-v7-3-canonical-parity}
(HEAD \texttt{7ee15a964bce51aba22a997e9bba31a389514ac2} at time of
writing). \\[0.5em]
\textbf{Authority rule applied throughout this Part}: where a
requirement below restates existing \texttt{NORMATIVE\_CANON.md} or
v7.2/v7.2.1 content, it is presented as normative (it already was).
Where a requirement has \emph{no} Canon/v7.2 precedent and derives
solely from the secondary implementation-target document, it was
originally marked \textbf{SOURCE-PENDING} --- offered as a working
assumption, not as confirmed original-author intent. \textbf{No content
in this Part amends \texttt{NORMATIVE\_CANON.md}.}
Any apparent tension with the Canon is resolved in the Canon's favor;
see \S\ref{sec:v73-typezero-disambiguation} and
\S\ref{sec:v73-ruin-disambiguation} for the two cases this build found.
\\[0.5em]
\textbf{2026-09-21 update}: following the closure of original-source
search (Human Sovereign instruction), 5 of the 6 SOURCE-PENDING items
were explicitly \textbf{ADOPTED} by Human Sovereign decision as
\emph{new, 2026 content grounded in surviving evidence} --- see
\texttt{docs/V7\_3\_ADOPTION\_RECORD.md}. \textbf{This is not a claim
that the original ``NRMO SYSTEM v7.3'' text has been located,
recovered, or reconstructed}; that text remains
\textbf{ORIGINAL SOURCE NOT AVAILABLE}, and each item marked
\textbf{ADOPTED} below is tagged with its adoption date and remains
distinguishable from Canon/v7.2-precedented text, per Canon \S10. One
item (\S\ref{sec:v73-strongengine-surface}, StrongEngine candidate-
surface categories) remains \textbf{SOURCE-PENDING and DEFERRED}; an
open-ended revision is proposed but not yet adopted. This build's
overall \textbf{PROVISIONAL} status (this notice, the frontmatter
notice, and the title page) is unchanged by this partial adoption.
}

\section{Relationship to DecisionCompass}
\label{sec:v73-decisioncompass-relationship}

Per the Founding Charter and \texttt{NORMATIVE\_CANON.md} \S11
(Conformance Rule), the authority order is:

\begin{center}
\texttt{NRMO normative specification} $\to$
\texttt{DecisionCompass implementation} $\to$
\texttt{conformance verification}
\end{center}

This Part is written to precede and govern DecisionCompass's v7.3
implementation, not to certify or adopt it. Where DecisionCompass's
implementation-scope document is used as evidence of what v7.3
requires, this Part treats that document the way any implementation
proposal is treated under \texttt{NORMATIVE\_CANON.md} \S9 and \S11:
informative, not self-authorizing. DecisionCompass's own conformance
matrix (\texttt{docs/nrmo\_v7\_3\_conformance.json}) is
implementation-only evidence (Table~\ref{tab:v73-classification-summary},
row category \texttt{IMPLEMENTATION-ONLY}) and is out of scope for this
Part beyond the cross-reference in
\S\ref{sec:v73-decisioncompass-conformance-summary}.

\chapter{Authority, Mode, and Lifecycle Model}
\label{ch:v73-authority-mode}

\section{Authority Hierarchy (restated, unchanged)}
\label{sec:v73-authority}

\begin{center}
\texttt{Human Sovereign} $\to$ \texttt{Vision (human-owned)} $\to$
\texttt{NRMO (governance/admissibility)} $\to$ \texttt{A\_allowed} $\to$
\texttt{StrongEngine \(\Omega\) Full (search/rank only)} $\to$
\texttt{recommendation} $\to$ \texttt{Human Sovereign (final)}
\end{center}

This is \texttt{NORMATIVE\_CANON.md} \S2 restated with the admissible
set made explicit as an intermediate stage; it introduces no new
authority claim. The strict invariant
$a_t \in A_t$ (Canon \S2.3, \S8) is unchanged.

\section{Four-Layer Mode/State Model}
\label{sec:v73-four-layer}

Four state axes must be represented and traced \emph{separately} and
must never be flattened into a single enum:

\begin{enumerate}
\item \textbf{Operational Mode}: \texttt{NORMAL | SAFE | VENTURE |
MISSION} --- Canon \S3.1, unchanged.
\item \textbf{Context Override}: \texttt{TRAINING(A--F,E+) | HARE |
NONE} --- Canon \S3.2--\S3.3, unchanged.
\item \textbf{Type ZERO Internal Gating State} --- see
\S\ref{sec:v73-typezero-disambiguation}; historical labels, disambiguated.
\item \textbf{Lifecycle State}: \texttt{ACTIVE | HOLD | EXIT |
SAFE\_EXIT} --- Canon \S3.4, unchanged.
\end{enumerate}

\noindent
The $A_{\mathrm{allowed}}$ extension \texttt{MISSION-DEFENSE =
active|inactive} (Canon \S3.6, Section~\ref{sec:mission-defense-c2}) is
an extension flag on axis (1)'s MISSION profile, not a fifth axis.

\subsection{Disambiguation: Type ZERO Internal Gating State vs Operational Mode}
\label{sec:v73-typezero-disambiguation}

\nrmobox{Resolved terminology-collision risk}{
\textbf{Finding}: the v7.3 implementation-target document specifies a
``Type ZERO Mode'' axis using the labels
\texttt{CORE|VENTURE|MISSION|SHUTDOWN}. Section~\ref{sec:type-zero-modes-c2}
already preserves exactly this 4-tuple as an explicit
\emph{historical, superseded} table --- superseded specifically
\emph{as the Operational Mode taxonomy} (\texttt{NORMATIVE\_CANON.md}
\S5). Reusing identical labels (\texttt{MISSION}, \texttt{VENTURE}) for
two different axes creates a serious readability and
implementation-drift risk: a trace value \texttt{MISSION} would be
ambiguous between ``Operational Mode = MISSION'' and ``Type ZERO
internal gate = MISSION'' without qualification. \\[0.5em]
\textbf{Resolution (Canon-priority reading, no Canon amendment)}:
Canon \S5 supersedes this label set only as the Operational Mode
taxonomy. It does not forbid Type ZERO from retaining an internal,
non-actuator (Canon \S6) gating-state vocabulary that happens to reuse
historical labels, \emph{provided it is never exposed or logged
without an explicit \texttt{type\_zero\_gate} qualifier distinct from
\texttt{operational\_mode}}. Any implementation or document that
prints a bare \texttt{MISSION} without stating which axis it belongs
to is non-conformant. This is a documentation/trace-schema obligation,
not a theory change.
}

Required trace discipline going forward: every occurrence of this
label set outside Section~\ref{sec:type-zero-modes-c2} itself must be
qualified as \texttt{type\_zero\_gate.\{CORE|VENTURE|MISSION|SHUTDOWN\}}
to prevent confusion with \texttt{operational\_mode.\{NORMAL|SAFE|VENTURE|MISSION\}}.

\section{HOLD State Machine}
\label{sec:v73-hold-state-machine}

\nrmobox{Correction recorded during self-audit (LOOP 6)}{
An earlier draft of this section classified the transition table below
as \texttt{ADDED / SOURCE-PENDING}, i.e.\ as having no direct v7.2
textual precedent. Re-reading Chapter~\ref{ch:theoretical-invariants}
(\S\ref{sec:inv-state-space}) during self-audit found that this exact
five-transition table already exists there, verbatim, as part of the
v2.0 theoretical invariants (\textbf{UNCHANGED} status in that
chapter's own terms). The classification below is corrected
accordingly. This correction is recorded rather than silently applied,
per the evidence-discipline requirement governing this Part.
}

\texttt{NORMATIVE\_CANON.md} \S3.4 establishes the lifecycle-state
family; Chapter~\ref{ch:theoretical-invariants} \S\ref{sec:inv-state-space}
already provides the transition table itself. The following table is
therefore \textbf{UNCHANGED} (restated for local readability in the
v7.3 context, not new content):

\begin{table}[ht]
\centering
\small
\begin{tabularx}{\linewidth}{lXX}
\toprule
\textbf{Transition} & \textbf{Trigger} & \textbf{Authority required} \\
\midrule
\texttt{ACTIVE $\to$ HOLD} & Human intervention required (e.g.\ empty
$A_{\mathrm{allowed}}$, materially insufficient information for a
major/irreversible candidate) & NRMO governance signal; no human
action needed to \emph{enter} HOLD \\
\texttt{ACTIVE $\to$ EXIT} & Planned termination (mission
success/exit criteria met) & Per the applicable MISSION exit criteria \\
\texttt{ACTIVE $\to$ SAFE\_EXIT} & Emergency termination & NRMO
governance signal (ruin-boundary or SHUTDOWN-equivalent trigger) \\
\texttt{HOLD $\to$ ACTIVE} & Resume & \textbf{Explicit Human Sovereign
approval only} (Canon \S2.1) \\
\texttt{HOLD $\to$ SAFE\_EXIT} & Escalation (HOLD condition
deteriorates or times out unresolved) & NRMO governance signal \\
\bottomrule
\end{tabularx}
\caption[HOLD state machine]{HOLD lifecycle transitions. UNCHANGED:
verbatim restatement of Chapter~\ref{ch:theoretical-invariants}
\S\ref{sec:inv-state-space} (v2.0 theoretical invariant, preserved
through v7.2), read together with Canon \S3.4.}
\label{tab:v73-hold-transitions}
\end{table}

\noindent
See also \S\ref{sec:hold-refinement} (\texttt{HOLD\_WITH\_OBSERVATION
| HOLD\_WITH\_SHADOW | HOLD\_WITH\_EXPIRY}) for the pre-existing
operational refinement of the single abstract \texttt{HOLD} state used
above; that refinement is orthogonal to this table and is unchanged by
it.

\paragraph{Empty admissible set (\textbf{ADOPTED --- v7.3, New Explicit
Adoption, Human Sovereign, 2026-09-21}).} If
$A_{\mathrm{allowed}}$ (equivalently
$A_t$) is empty, the execution layer returns \texttt{HOLD} and
generates \emph{no} execution candidate. This follows directly from
Canon \S8 (execution may never expand $A_t$ on its own authority): if
governance has narrowed $A_t$ to $\emptyset$, no admissible candidate
can exist by construction, and any output claiming otherwise would
violate $a_t \in A_t$. Adopted by explicit Human Sovereign decision,
2026-09-21 (see \texttt{docs/V7\_3\_ADOPTION\_RECORD.md}, item 1); this
is new-explicit v7.3 content, not a claim that the unlocated original
``NRMO SYSTEM v7.3'' text stated this rule.

\chapter{Norn: Observation Responsibilities}
\label{ch:v73-norn}

\texttt{NORMATIVE\_CANON.md} \S7.1 already establishes canonical Norn
as an interpretive governance perspective that is \emph{not} the
sovereign, not a decision-maker, not StrongEngine, not an execution
engine. The following itemized responsibility list is
\textbf{ADOPTED (v7.3, New Explicit Adoption --- Human Sovereign,
2026-09-21)}: an elaboration consistent with \S7.1, not previously
itemized at this granularity in Canon/v7.2 text. See
\texttt{docs/V7\_3\_ADOPTION\_RECORD.md}, item 2; this is new-explicit
content, not restored original text.

\section{Required Observable Responsibilities}

\begin{itemize}
\item Branch record: proposed / admissible / selected / rejected
alternatives.
\item Procedure audit: whether the ordered evaluation of
\S\ref{sec:v73-decision-procedure} was actually followed by the
runtime, including detection of missing or reordered steps.
\item Conceptual-drift observation, using the descriptive
classification \texttt{DRIFT | SIGNAL | WITHIN\_OWN\_RANGE |
INSUFFICIENT\_HISTORY} (\textbf{ADOPTED --- v7.3, New Explicit
Adoption, Human Sovereign, 2026-09-21; see
\texttt{docs/V7\_3\_ADOPTION\_RECORD.md}, item 3}; descriptive only;
\textbf{no Norn-owned decision threshold is permitted} --- ``audit has
no scale'', directly entailed by Canon \S7.1 + \S8).
\item Passive observation: stagnation, opportunity loss, adaptation
delay.
\item Post-decision, one-way observation that cannot alter the
decision already made.
\item Upstream observation immediately after normalized World
Model/Observation, where the runtime makes this feasible.
\end{itemize}

\section{Explicit Prohibitions (restated from Canon, made explicit)}
\label{sec:v73-norn-prohibitions}

Canonical Norn \textbf{must not}: veto; execute; set thresholds;
choose the route, Vision, or Mission; or return the final conclusion.
Each of these is already entailed by Canon \S7.1 (``not the final
decision maker'', ``not an execution engine'') read together with
\S8 (governance/execution separation is exhaustive: Norn is neither
governance nor execution). This section is \textbf{CLARIFIED}, not new
authority.

\chapter{Decision Procedure}
\label{ch:v73-decision-procedure}

\section{Ordered Evaluation}
\label{sec:v73-decision-procedure}

The following ordered procedure is \textbf{ADOPTED (v7.3, New Explicit
Adoption --- Human Sovereign, 2026-09-21)} as a single composite
sequence; see \texttt{docs/V7\_3\_ADOPTION\_RECORD.md}, item 4. Each
individual stage already has independent Canon or chapter grounding
(cited inline); the \emph{ordering itself} and its packaging as one
traceable procedure had no prior single-document Canon/v7.2 precedent
and is not asserted as pre-existing Canon text --- only as consistent
with it, and now adopted as new-explicit v7.3 content.

\begin{enumerate}
\item Vision / Mission alignment (Canon \S2.1, \S2.4; MISSION
definition, \S4/Canon).
\item Ruin boundary check (Chapter~\ref{ch:theoretical-invariants};
Chapter~\ref{ch:time-horizon-layer}; \S\ref{sec:v73-ruin-disambiguation}
below).
\item Reversibility check (R1-FIX Boundary Contract,
\S\ref{sec:r1fix-boundary-contract}).
\item $A_{\mathrm{allowed}}$ narrowing
(Chapter~\ref{ch:a-allowed}, Chapter~\ref{ch:a-allowed-dictionary}).
\item Information sufficiency check (major/irreversible lock-in is
blocked under material information insufficiency ---
\S\ref{sec:v73-a-allowed-refinement}).
\item Execution candidate generation
(\S\ref{sec:v73-strongengine-surface}).
\item Observation indicator recording (Norn,
\S\ref{ch:v73-norn}).
\item Exit condition statement.
\item Reevaluation condition/date statement.
\item Return final authority to Human Sovereign (Canon \S2.1).
\end{enumerate}

\noindent
Norn audits whether this order was followed; Norn cannot alter it
(\S\ref{sec:v73-norn-prohibitions}).

\chapter{Ruin Taxonomy}
\label{ch:v73-ruin-taxonomy}

\section{Disambiguation: Two Non-Identical Ruin Axes}
\label{sec:v73-ruin-disambiguation}

\nrmobox{Resolved terminology risk --- do not read as one 4-member enum}{
\textbf{Finding}: the v7.3 implementation-target document lists ``Hard
Ruin, Soft Ruin, Active Ruin, Passive Ruin'' as four items to keep
distinct in trace fields. Read literally as one flat taxonomy, this
conflates two chapters with \emph{different scope and different
axis-meaning}:
\begin{itemize}
\item \textbf{Hard Ruin / Soft Ruin} (Chapter~\ref{ch:investment-sop},
\S47--51): a \emph{severity} axis, defined specifically within the
Investment SOP domain. Hard Ruin = unrecoverable (bankruptcy,
permanent capital loss); Soft Ruin = recoverable-in-principle
strategic failure (forced Risk-Off transition).
\item \textbf{Active Ruin / Passive Ruin}
(\S\ref{sec:omega-active-vs-passive}): a \emph{causal-origin} axis,
defined domain-generally in the StrongEngine $\Omega$ Full chapter.
Active Ruin = direct collapse from reckless action; Passive Ruin =
erosion of optionality from avoidance/rigidity/adaptation failure.
\end{itemize}
These are not four siblings of one enum; they are two independent
axes from two different chapters, one domain-scoped and one
domain-general, that may co-occur (e.g.\ a Soft Ruin in the investment
domain may have either an Active or a Passive causal origin).
}

\textbf{Resolution}: trace schemas that need both axes must record
them as \texttt{ruin\_severity \(\in\) \{Hard, Soft\}} (Investment
domain only, may be \texttt{N/A} elsewhere) and
\texttt{ruin\_origin \(\in\) \{Active, Passive\}} (domain-general)
as two separate fields, never as one four-valued field. This
resolves the CONFLICT identified in the differential table
(\texttt{docs/V7\_3\_DIFFERENTIAL\_TABLE.md}, row 6.1) to
\textbf{CLARIFIED} without requiring any Canon change (Canon does not
mention any of the four terms; the ambiguity was entirely
chapter-to-chapter).

\section{Passive Ruin: Avoidability-Window Definition (v7.2.1, adopted)}
\label{sec:v73-passive-ruin-window}

This is the single strongest-evidence item in this Part: the
avoidability-window definition of Passive Ruin was already adopted in
v7.2.1, \emph{before} any v7.3 implementation work began
(\texttt{docs/v72\_1\_validation\_record.md},
\texttt{docs/v72\_implementation\_integration\_map.md} \S2--3). It is
formally carried into v7.3 as \textbf{CLARIFIED} (re-affirmed, not
newly introduced):

\begin{itemize}
\item \textbf{Leading indicator} (primary): at the moment of choice,
does the candidate action close an avoidability window that was open
in the state immediately before the choice? (\texttt{window\_open()},
\texttt{closes\_window()}, diff-in-diff against the do-nothing
baseline --- \texttt{docs/v72\_implementation\_integration\_map.md}
\S2--3.)
\item \textbf{Lagging indicator} (retained, not deleted): the legacy
streak-based signal (time spent in a degraded-state region) is kept
as a secondary, after-the-fact diagnostic, exactly as
v7.2.1 specified. It is not deleted per the ``preserve, do not erase''
guardrail (\texttt{NORMATIVE\_CANON\_ADOPTION.md} \S Guardrails).
\end{itemize}

\noindent
\textbf{Passive Ruin must not be redefined as mere inactivity.} An
implementation that flags Passive Ruin purely from an inactivity
streak, without an observed avoidability-window closure, is
non-conformant with this definition (this restates
\texttt{docs/v72\_1\_validation\_record.md} \S3--4, not a new rule).

\chapter{$A_{\mathrm{allowed}}$ and StrongEngine Refinement}
\label{ch:v73-aallowed-strongengine}

\section{Admissibility Refinement}
\label{sec:v73-a-allowed-refinement}

Building on Chapter~\ref{ch:a-allowed} and
Chapter~\ref{ch:a-allowed-dictionary} (\textbf{UNCHANGED} baseline), an
action is inadmissible (excluded from $A_{\mathrm{allowed}}$) when it:

\begin{enumerate}
\item creates direct irreversible loss beyond the applicable
reversibility floor (R1-FIX Boundary Contract,
\S\ref{sec:r1fix-boundary-contract});
\item locks in a major/high-commitment decision under material
information insufficiency;
\item removes a required withdrawal or review capability; or
\item violates a sovereignty constraint (i.e.\ would remove or bypass
the Human Sovereign's ability to accept, reject, defer, or terminate
--- Canon \S2.1).
\end{enumerate}

This is \textbf{CLARIFIED}: each clause is independently grounded in
existing chapters; their joint statement as an admissibility checklist
is new packaging, not new substance.

\section{StrongEngine Candidate Surface}
\label{sec:v73-strongengine-surface}

Within $A_{\mathrm{allowed}}$, Chapter~\ref{ch:strong-engine} and
Chapter~\ref{ch:omega-full-integrated} already establish that StrongEngine
$\Omega$ Full searches and ranks candidates without redefining the
ruin boundary. The following candidate-category checklist is
\textbf{ADDED / SOURCE-PENDING} (a specific, non-exhaustive minimum
surface, not previously enumerated at this granularity):
safe/conservative; standard; assertive/forward; information-acquisition
probe; bounded sequence/checkpoint plan; costed HOLD (where
$A_{\mathrm{allowed}}$ admits it); exit/recovery (where required).

\noindent
\textbf{Invariant (UNCHANGED)}: no advisory or presentation layer may
substitute an inadmissible action for the StrongEngine's selection.
This is the strict invariant $a_t \in A_t$ (Canon \S2.3, \S8) restated,
not a new rule.

\chapter{Trace and DecisionRecord Schema}
\label{ch:v73-trace-schema}

\texttt{NORMATIVE\_CANON.md} \S10 requires that publication version,
baseline version, component version, and implementation-snapshot
version remain distinguishable, but does not itself enumerate a
decision-trace schema. The following field list is
\textbf{ADOPTED (v7.3, New Explicit Adoption --- Human Sovereign,
2026-09-21)}, offered as the minimum information needed to audit a
single decision against every requirement in this Part; see
\texttt{docs/V7\_3\_ADOPTION\_RECORD.md}, item 6:

\begin{itemize}
\item authority/spec version (e.g.\ \texttt{NRMO v7.3} tag) and
component versions;
\item operational mode / context override / Type ZERO internal gate
(qualified per \S\ref{sec:v73-typezero-disambiguation}) / lifecycle
state / MISSION-DEFENSE flag;
\item Vision/Mission reference, where supplied by the Human Sovereign;
\item ruin results: \texttt{ruin\_severity} and \texttt{ruin\_origin}
as two separate fields (\S\ref{sec:v73-ruin-disambiguation});
\item reversibility classification;
\item admissible and vetoed candidates, with reasons;
\item selected action and StrongEngine score metadata;
\item observation indicators (Norn);
\item exit conditions;
\item reevaluation conditions/date;
\item Norn branch / procedure / drift observations
(\S\ref{ch:v73-norn});
\item Human Sovereign's recorded acceptance, modification, or
rejection, where captured.
\end{itemize}

\chapter{Summary Classification and DecisionCompass Cross-Reference}
\label{ch:v73-summary}

\section{Classification Summary}

\begin{table}[ht]
\centering
\small
\begin{tabularx}{\linewidth}{lXr}
\toprule
\textbf{Classification} & \textbf{Meaning} & \textbf{Rows (differential table)} \\
\midrule
UNCHANGED & Restates existing Canon/v7.2 with no delta & 14 \\
CLARIFIED & Adds precision consistent with Canon/v7.2 & 7 \\
ADOPTED (v7.3, New Explicit Adoption --- Human Sovereign, 2026-09-21) & New content, explicitly adopted on surviving evidence; not a claim of Canon/v7.2 precedent or of restored original text & 5 \\
ADDED (SOURCE-PENDING, DEFERRED) & New, uncontradicted by Canon/v7.2; adoption decision pending (row 10.1 only) & 1 \\
SUPERSEDED & v7.3 explicitly replaces prior content & 0 \\
CONFLICT (resolved) & Apparent contradiction, resolved without Canon change & 1 \\
IMPLEMENTATION-ONLY & DecisionCompass runtime/CI concern & 2 \\
HISTORICAL & Refers to already-non-normative material & 1 \\
\bottomrule
\end{tabularx}
\caption[v7.3 differential classification summary]{Summary of the
full row-by-row reconciliation in
\texttt{NRMOIntegrated/docs/V7\_3\_DIFFERENTIAL\_TABLE.md}, updated per
the 2026-09-21 Human Sovereign adoption decision recorded in
\texttt{docs/V7\_3\_ADOPTION\_RECORD.md}. No row required a
\texttt{NORMATIVE\_CANON.md} change. Row 3.1 (HOLD
transition table) was reclassified from ADDED/SOURCE-PENDING to
UNCHANGED during LOOP 6 self-audit after locating verbatim precedent
in Chapter~\ref{ch:theoretical-invariants} \S\ref{sec:inv-state-space};
see that section's note for the correction record.}
\label{tab:v73-classification-summary}
\end{table}

\section{DecisionCompass Conformance Cross-Reference}
\label{sec:v73-decisioncompass-conformance-summary}

As of DecisionCompass branch \texttt{agent/nrmo-v7-3-canonical-parity}
(HEAD \texttt{7ee15a964bce51aba22a997e9bba31a389514ac2}), the
implementation's own conformance matrix
(\texttt{docs/nrmo\_v7\_3\_conformance.json}) records all tracked
requirements as \texttt{UNIT\_VERIFIED}, \texttt{SCENARIO\_VERIFIED},
or \texttt{CI\_VERIFIED}, with one item
(\texttt{physical-device}) recorded as
\texttt{NOT\_RUN\_NO\_CONNECTED\_DEVICE} (conditional on hardware
availability, not a specification gap). This cross-reference is
informational: it describes the implementation's self-reported state
at a pinned commit and does \textbf{not} certify conformance on behalf
of NRMO, and does not make the implementation the source of this
Part's content (\texttt{NORMATIVE\_CANON.md} \S11).

\section{Non-Goals of This Part}

Restating the implementation-target document's own non-goals, which
are themselves consistent with \texttt{NORMATIVE\_CANON.md} \S9 and
\S11 (\textbf{UNCHANGED}):

\begin{itemize}
\item Experimental prior work is not adopted merely because it exists.
\item This Part does not require any particular service-extraction or
deployment architecture.
\item Historical implementation-specific task routers (e.g.\ Shinobi)
must not be conflated with canonical Norn (Canon \S7.2).
\item \textbf{NRMO theory is not rewritten for implementation
convenience} --- where this build found tension between an
implementation-target document and the Canon, the Canon prevailed
(\S\ref{sec:v73-typezero-disambiguation},
\S\ref{sec:v73-ruin-disambiguation}), and no Canon text was altered.
\end{itemize}
